\documentclass[12pt]{uz-exame} 
\usepackage{bbb,amssymb,amsfonts,amsmath}
\newcommand{\rem}[1]{}
\begin{document} 
\examtitle{BSc Honours in Mathematics 1}
\examdate{May/June 2013}
\hours{2}
\lastA{5}
\rubric{\science}
\lastpage{4}
\coursename{CALCULUS 2}
\coursenumber{HMTH111}
%\setcounter{exsec}{17}
\makeheading
\sectionAscience
\begin{question}
Define what is meant by saying that 
\begin{itemize}
\item[(a)] a function $f(x,y)$ has a \textbf{limit}, $L$, at a point $(a,b)$ as $(x,y)$ approaches $(a,b)$. \marks{2}
\item[(b)] a function $f(x,y)$ has a \textbf{partial derivative} $f_x(a,b)$ with respect to $x$ at a point $(a,b)$. \marks{2}
\item[(c)] the series $\displaystyle \sum _{n=1}^\infty a_n$ is \textbf{absolutely convergent}.\marks{2}
\item[(d)] the series $\displaystyle \sum _{n=1}^\infty a_n$ is \textbf{conditionally convergent}.\marks{2}
\end{itemize}
\end{question}
\begin{question}
\begin{itemize}
\item[(a)] Find the area of the triangle with vertices $(-1,2,3), (2,-3,1)$ and $(1,-1,2)$.\marks{4}
\item[(b)] Find the area of the parallelogram determined by ${\bf{a}}=3{\bf{i}}+4{\bf{j}}$ and ${\bf{b}}={\bf{i}}+{\bf{j}}+{\bf{k}}$.\marks{4}
\end{itemize}
\end{question}
\begin{question}
 Let $f(x,y)$ be a function of $x$ and $y$ where $x=e^{u-v}\cos (u+v)$ and $y=e^{u-v}\sin (u+v)$. Show that 
\[\dfrac{\partial f}{\partial u}+\dfrac{\partial f}{\partial v}=2x\dfrac{\partial f}{\partial y}-2y\dfrac{\partial f}{\partial x}.\]\marks{6}
\end{question}
\begin{question}
Let \[f(x,y) = \left\{ 
\begin{array}{l l}
  \dfrac{xy}{x^2+y^2}, & \quad (x,y)\neq (0,0)\\
  0, & \quad (x,y)=(0,0).\\ \end{array} \right. \]
  \begin{itemize}
  \item[(a)] Determine whether or not $\dfrac{\partial f}{\partial x}$ and $\dfrac{\partial f}{\partial y}$ exists at the origin, and find their values if they exist.\marks{4}
  \item[(b)] Does $\displaystyle \lim _{(x,y)\rightarrow (0,0)} f(x,y)$ exist? \marks{4}
  \end{itemize}
\end{question}
\begin{question}
\begin{itemize}
\item[(a)] Determine whether or not the series converges and find its sum if it converges
\[\sum _{r=1}^\infty \dfrac{20}{(7r-3)(7r+4)}=\dfrac{20}{4\cdot 11}+\dfrac{20}{11\cdot 18}+\dfrac{20}{18\cdot 25}+\cdots +\dfrac{20}{(7n-3)(7n+4)} +\cdots .\]\marks{5}
\item[(b)] By interchanging the order of integration or otherwise, evaluate the integral \[ \int _{0}^1\int _{3y}^3 e^{x^2}\, dxdy.\]\marks{5}
\end{itemize}
\end{question}
\sectionBscience
\begin{question}
\begin{itemize}
\item[(a)] \begin{itemize}
\item[(i)] Write down the equation of the tangent plane to the ellipsoid \\ $2x^2+4y^2+z^2-45=0$ at the point $(2,-3,-1)$. \marks{5}
\item[(ii)] Find the equation of the plane passing through the points\\ $(2,-1,2), \, (-1,2,-3)$ and $(4,1,0)$. \marks{5}
\end{itemize}
\item[(b)]
\begin{itemize}
\item[(i)] Find the angle between the vectors ${\bf{a}}=2{\bf{i}}+4{\bf{j}}+6{\bf{k}}$ and ${\bf{b}}={\bf{i}}-3{\bf{j}}+2{\bf{k}}$. \marks{4}
\item[(ii)] Find the parametric equations of the line of intersection of the planes 
\[2x+3y-z=-3 \quad \mbox{and} \quad 4x+5y+z=1\] and calculate the angle between them. \marks{8}
\end{itemize}
\item[(c)] \begin{itemize}
\item[(i)] Find the parametric and symmetric equations for the straight line through $(2,5,-7)$ and $(4,3,8)$. \marks{4}
\item[(ii)] Find the distance between the point $(5,12,-13)$ and the plane\\ $3x+4y+5z=12$. \marks{4}
\end{itemize}
\end{itemize}
\end{question}
\begin{question}
\begin{itemize}
\item[(a)] Using the definition of a limit, show that $\displaystyle \lim _{(x,y)\rightarrow (2,1)}(xy-3x+4)=0$. \marks{6}
\item[(b)] Find and classify the stationary points of the function \[f(x,y)=x^3+y^3-3x-12y+20.\] \marks{6}
\item[(c)] Parabolic co-ordinates $(u,v)$ are defined implicitly in terms of the Cartesian co-ordinates $(x,y)$ by the pair of equations
\[x=\dfrac{u^2-v^2}{2} \quad \mbox{and}\quad y=uv.\]
\begin{itemize}
\item[(i)] Obtain expressions for $\dfrac{\partial u}{\partial x},\, \dfrac{\partial u}{\partial y} , \, \dfrac{\partial v}{\partial x}$ and $\dfrac{\partial v}{\partial y}$. \marks{8}
\item[(ii)] Given that $f(x,y)=\phi (u,v)$. Obtain expressions for $\dfrac{\partial f}{\partial x}$ and $\dfrac{\partial f}{\partial y}$ in terms of $\dfrac{\partial \phi}{\partial u}, \, \dfrac{\partial \phi}{\partial v}, \, u$ and $v$ and deduce that  \[\left(\dfrac{\partial f}{\partial x}\right)^2+\left(\dfrac{\partial f}{\partial y}\right)^2=\dfrac{1}{u^2+v^2}\left[\left(\dfrac{\partial \phi}{\partial u}\right)^2+\left(\dfrac{\partial \phi}{\partial v}\right)^2\right].\]\marks{10}
\end{itemize}
\end{itemize}
\end{question}
\begin{question}
\begin{itemize}
\item[(a)] Sketch the region of integration and evaluate the following integrals.
\begin{itemize}
\item[(i)] $\displaystyle \int _0^1 \int _{y}^{\sqrt{y}} \, dxdy$.
\item[(ii)] $\displaystyle \iint \limits _D \sqrt{x^2+y^2}\, dydx$ where $D$ is the region bounded by $x^2+y^2\geq 4$ and\\ $x^2+y^2\leq 9$. \marks{3,3}
\end{itemize}
\item[(b)] \begin{itemize}
\item[(i)] Evaluate $\displaystyle \int _0^1\int _0^1 \int _{\sqrt{x^2+y^2}}^2 xyz\, dzdydx.$ \marks{5}
\item[(ii)] Evaluate the following integral, using an appropriate transformation or otherwise \[\int _0^1 \int _0^{\sqrt{1-y^2}} \dfrac{dxdy}{1+x^2+y^2}.\] \marks{5}
\end{itemize}
\item[(c)] Let $u=x^2-y^2$ and $v=2xy$.
\begin{itemize}
\item[(i)] Show that $(x^2+y^2)^2=u^2+v^2$. \marks{3}
\item[(ii)] Hence, find $\dfrac{\partial (x,y)}{\partial (u,v)}$ in terms of $u$ and $v$. \marks{5}
\item[(iii)] Hence, evaluate $\displaystyle \iint \limits _R (x^2+y^2)\, dxdy$ where $R$ is the region in the first quadrant of the $xy-$plane bounded by the curves $x^2-y^2=1, \, x^2-y^2=9, \, xy=2$ and $xy=4$. \marks{6}
\end{itemize}
\end{itemize}
\end{question}
\begin{question}
\begin{itemize}
\item[(a)] Test for convergence or divergence for each of the following infinite series by applying an appropriate test.\\
(i)\, $\displaystyle \sum _{n=1}^\infty \dfrac{n}{2^n}$ \quad (ii)\, $\displaystyle \sum _{n=1}^\infty \dfrac{n-7}{4n^3+3}$ \quad (iii)\, $\displaystyle \sum _{n=1}^\infty \dfrac{n}{\sqrt{n^4+1}}$. \marks{3,3,3}
\item[(b)] \begin{itemize}
\item[(i)] Find the sum of the series $\displaystyle \sum _{n=1}^\infty\dfrac{1}{n(n+1)}$. \marks{4}
\item[(ii)] Consider the series $\displaystyle \sum _{n=1}^\infty ar^{n-1}$ where $a$ and $r$ are constants. Determine when the series converges and when it diverges and find the sum when it converges. \marks{5}
\end{itemize}
\item[(c)] Test the following series for absolute and conditional convergence and determine the interval of convergence.
\begin{itemize}
\item[(i)] $\displaystyle \sum _{n=1}^\infty \dfrac{(-1)^n x^n}{n2^n}$. 
\item[(ii)] $\displaystyle \sum _{n=1}^\infty \dfrac{n^3x^n}{n!} $. 
\end{itemize}\marks{6,6}
\end{itemize}
\end{question}
\endexam
\end{document} 